\subsection{Verificación de informes durables}\label{sec:pruebas-informes}

La aceptación distingue autorización, consistencia de captura, persistencia,
renderizado y presentación. Una prueba de generación aislada no demuestra que
un archivo sólo pueda descargarse por su solicitante vigente ni que ambos
formatos sobrevivan a reinicios y restauraciones. Tampoco una página del
selector constituye un conteo global. Los grupos focales se registran por
su alcance y no se suman como si fueran una regresión completa del producto.

En la comprobación local de composición y supervisión aprobaron 43 pruebas.
El grupo focal de renderizado y paginación aprobó 15; el selector de litigantes
aprobó seis en aplicación y cuatro con PostgreSQL aislado. El contrato HTTP de
informes aprobó 14 y el cliente de interfaz, 19 verificaciones Node. El selector
incluye un colega compartido únicamente en expedientes cerrados, paginación de
litigantes activos sin expedientes para Owner, retirada de pertenencia y rechazo
de roles ajenos incluso para una página vacía. La aplicación rechaza además
respuestas incoherentes de alcance, orden, cursor o reloj y vuelve a autenticar
al solicitante antes de entregar el resultado.

La paginación del PDF debe conservar texto y filas en los saltos de página, sin
recortar líneas ni dejar espacio reservado que fuerce una página innecesaria.
La comprobación de recursos debe combinar cantidad de casos, asignaciones y
carga con textos representativos y de anchura máxima. Un resultado dentro del
máximo de filas aún puede exceder páginas, bytes o el presupuesto del proceso.
Por ello, los valores de la \cref{sec:impl-informes} son límites de admisión y
ejecución, no una medición de rendimiento ni garantía de aceptar cualquier
contenido con esa cantidad de filas.

En el navegador con HTTP controlado se acreditaron trece escenarios distintos;
las repeticiones focales no se cuentan como escenarios adicionales. Con servicios
reales aprobaron dos recorridos de informes, uno Owner y otro Litigante, en
10.6 y 7.6 segundos respectivamente. El ejecutor de navegador duró 23.4 segundos;
la campaña completa, 207.863 segundos. Ambos recorrieron solicitud, actualización,
descarga exacta de PDF y CSV, aviso, navegación y nueva autenticación, conservando
los artefactos. La inspección de cuatro capturas, a 1440 y 390 píxeles para ambos
roles, aprobó la presentación de escritorio y móvil.

Dos pruebas focales de reutilización del plan tipográfico conservaron exactamente
los identificadores de glifo y sus posiciones, incluidos texto largo, acentos y
ligaduras. La captura representativa de sesenta expedientes produjo el mismo PDF
de veinticinco páginas, byte por byte, en 1.6041 segundos frente a los 2.7404
segundos previos, sin modificar los límites. La inspección visual de sus páginas
1, 2, 3, 13 y 25 aprobó.

La sustitución posterior de las bibliotecas tipográficas por HarfRust y Skrifa,
con mantenimiento vigente, conservó ese PDF de 1654401 bytes exactamente y
produjo una observación aislada de 0.566156 segundos. La compilación previa
consumió 59.44 segundos y se contabilizó aparte. Son muestras individuales del
renderizador; no acreditan una reducción equivalente del CI. La revisión visual
del PDF exportado se conserva por su identidad de bytes.

En la medición anterior a esa sustitución, la captura de mil expedientes,
diez mil asignaciones y mil entradas de carga aprobó CSV en 0.382 segundos. Su PDF continuó devolviendo
\texttt{CapacityExceeded}: 15.197 segundos en la medición posterior y 15.284 en
la anterior. La categoría pública no identifica cuál de los presupuestos se
alcanzó; no se atribuye el rechazo sólo a CPU ni se declara generación del PDF
máximo. Son observaciones del renderizador aislado, no una frontera exhaustiva
de capacidad ni una aceptación de producción.

La aceptación API integrada aprobó con salida cero en 326.299 segundos. El
recorrido HTTP creó tres trabajos y comprobó la repetición de una misma
operación, los pares PDF/CSV, sus identidades y digests, los avisos propios y el
primer acuse. En dos trabajos se observaron cola, procesamiento en renderizado
y disponibilidad; el trabajo después revocado se observó en cola y disponible.
El sondeo no acredita haber observado todos los estados intermedios.

Los reinicios mediante \texttt{SIGTERM} y \texttt{SIGINT} concluyeron en cinco
segundos cada uno y conservaron el histórico exacto. La restauración de
PostgreSQL y la renovación de MFA conservaron los bytes de PDF y CSV, sus
identidades, los avisos leídos y no leídos y el rechazo completo de la captura
revocada. Retirar la pertenencia a uno de los expedientes capturados impidió
descargar el informe entero; una operación idempotente no permitió recuperar
el archivo revocado. Los roles Client y Paralegal fueron rechazados.
Estos resultados están separados de las pruebas focales y de la regresión
global en \rutarepo{docs/verification-report.md}.

El criterio de la \cref{tab:rf-19} exige comprobar una interfaz responsiva
cuando la generación supera tres segundos, conservando la sesión y la consulta
de progreso. La notificaci\'on interna y la duraci\'on observada no demuestran
por s\'i solas el correo ni la predicci\'on inicial de la \cref{tab:cu-16}; las
pruebas de esos mecanismos se distinguen a continuaci\'on. Tampoco la carga de
expedientes acredita por s\'i sola la actividad atribuida a cada cuenta ni los
dem\'as tipos de reporte.

\paragraph{Estimaci\'on y disponibilidad por correo.}
Las pruebas incorporadas contrastan el m\'inimo de historial comparable, el
escalamiento por volumen, la mediana y el umbral estricto de dos segundos.
La persistencia comprueba la predicci\'on inicial al recuperar una operaci\'on,
la espera de la cola y la publicaci\'on at\'omica de intenci\'on, artefactos, aviso
y auditor\'ia. El consumidor comprueba identidad y acceso antes del transporte,
rechaza concesiones y ventanas vencidas, descarta resultados atrasados y conserva mensaje
e identidad de entrega despu\'es de una respuesta incierta o un reinicio.
Las pruebas de restauraci\'on y de inventario se mantienen separadas del
transporte capturado; ninguna acredita recepci\'on en una bandeja externa.

Las comprobaciones de HTTP y Qadra distinguen la estimaci\'on inicial de la
duraci\'on observada, admiten ausencia de datos en respuestas anteriores y
rechazan valores num\'ericos o estados incompatibles. Los recorridos definidos
para escritorio y m\'ovil muestran el estado pendiente, la aceptaci\'on del
proveedor, el rechazo, la incertidumbre y la retirada de acceso sin alterar
las descargas. La configuraci\'on y supervisi\'on comprueban desactivaci\'on por
defecto, canal expl\'icito, pausas entre intentos y parada conjunta. Sus resultados
se registran por ejecuci\'on en \rutarepo{docs/verification-report.md}; no se
suman a los conteos hist\'oricos anteriores.

La aceptaci\'on local con servidor real aprob\'o dos recorridos de navegador,
Owner y Litigante, de 15.1 y 8.2 segundos respectivamente. Playwright inform\'o
28.7 segundos y el comando completo, incluida preparaci\'on, 329.708 segundos.
La inspecci\'on de cuatro capturas de escritorio y m\'ovil aprob\'o la presentaci\'on.
Esta ejecuci\'on conserva el recorrido autorizado de los informes y se registra
separada de las mediciones hist\'oricas anteriores.

La aceptaci\'on nativa de disponibilidad aprob\'o en 35.992 segundos mediante HTTP,
PostgreSQL, Redis, contrase\~na Argon2 y MFA. Con dos predecesores en cola, el
informe conserv\'o su estimaci\'on inicial de 3360 milisegundos y alcanz\'o
disponibilidad en 3.453 segundos. El detalle pas\'o de solicitud en cola sin
intenci\'on de correo a disponible con aviso pendiente, y luego a aceptaci\'on
registrada por el transporte capturado. La nueva autenticaci\'on mediante un
c\'odigo de recuperaci\'on conserv\'o los bytes exactos de PDF y CSV; la cadena de
auditor\'ia permaneci\'o \'integra. Este ensayo comprueba el aviso elegible y la
recuperaci\'on del informe, sin acreditar recepci\'on externa ni activar correo
en producci\'on. Los controles de cierre e integraci\'on no se deducen de estas
aceptaciones focales y se distinguen al final de esta validaci\'on.

La validaci\'on independiente conserv\'o el estimador fijado y complet\'o primero
18 observaciones de calibraci\'on: tres por modalidad y volumen de 3, 80 y 900
expedientes. Despu\'es ejecut\'o 28 solicitudes nuevas: cuatro por modalidad y
volumen, m\'as una con predecesor en cola y otra con fallo transitorio por
modalidad, ambas de 80 expedientes. Cada solicitud utiliz\'o una copia nueva de
la base de calibraci\'on correspondiente, cuyo historial permaneci\'o congelado;
las duraciones de validaci\'on no se incorporaron a ese historial.

El criterio exige que al menos el 90\% de las solicitudes con predicci\'on tenga
un error absoluto no mayor que el m\'aximo entre dos segundos y el 30\% de su
duraci\'on real desde solicitud hasta disponibilidad. Las 28 solicitudes tuvieron
estimaci\'on inicial y alcanzaron disponibilidad: cobertura del 100\%, sin
solicitudes excluidas por falta de estimaci\'on. Cumplieron 26 de 28, equivalentes
al 92.857\%, como resume la \cref{tab:precision-estimacion-informes}. Las dos
solicitudes fuera de tolerancia permanecen en el denominador.

\begin{table}[htbp]
  \centering
  \caption{Validaci\'on independiente de la estimaci\'on inicial de informes.}
  \label{tab:precision-estimacion-informes}
  \begin{tabularx}{\linewidth}{@{}Xrrr@{}}
    \toprule
    Modalidad & Solicitudes & Predicciones & En tolerancia \\
    \midrule
    Estado y carga & 14 & 14 & 13 \\
    Actividad registrada & 14 & 14 & 13 \\
    \midrule
    Total & 28 & 28 & 26 \\
    \bottomrule
  \end{tabularx}
\end{table}

Los 24 casos individuales y los dos de cola cumplieron la tolerancia. Los dos
casos de reintento recibieron una falla \texttt{RenderUnavailable} inyectada y
respetaron la espera programada de treinta segundos antes de volver a usar el
renderizador aislado. Alcanzaron disponibilidad en 31.704 segundos para estado
y 31.475 para actividad, frente a predicciones iniciales de 1.690 y 1.461
segundos, respectivamente. El incidente imprevisto explica el desv\'io observado;
no se retir\'o del resultado ni se reemplaz\'o la predicci\'on inicial.

La ejecuci\'on nativa aprob\'o en 321.402 segundos con PostgreSQL 18.6 sobre Linux
x86\_64, un solo hilo de pruebas y temporales en btrfs, sin tr\'afico de producci\'on.
Esa duraci\'on incluye calibraci\'on y preparaci\'on, y no es latencia por solicitud.
La modalidad de actividad utiliz\'o exclusivamente cargas documentales originales,
sin actuaciones procesales ni plazos atendidos; el ensayo no valida el tiempo
de todo el cat\'alogo de actividad ni garantiza precisi\'on en otros contenidos,
equipos o cargas. Sus observaciones no sustituyen las mediciones hist\'oricas del
renderizador ni se mezclan con ellas. La integraci\'on, el despliegue y la
habilitaci\'on del canal con Resend en la instalaci\'on productiva, registrados en
\rutarepo{docs/technical-closure.md}, y la comprobaci\'on operativa del correo se
distinguen de esta aceptaci\'on con transporte de prueba: en la aceptaci\'on
productiva del 12 de octubre de 2026, registrada en
\rutarepo{docs/production-acceptance.md}, los dos informes observados mostraron
estimaci\'on no disponible y duraci\'on menor de un segundo, sin ejercitar la condici\'on de aviso por correo.

Estas comprobaciones no sustituyen la regresión global ni la validación de
usabilidad con participantes del despacho. La
\cref{anexo:informes} identifica las fuentes reproducibles; su existencia no
acredita ejecución, aprobación o despliegue.

\paragraph{Actividad atribuida por cuenta y periodo.}
El informe de actividad contrasta tres cantidades contra los registros originales,
usando el instante de operación y no la fecha judicial declarada. La comprobación
focal con PostgreSQL 16.15 aislado aprobó diez casos: cargas, hechos, actos de
recursos, sesiones/resultados, decisiones cautelares y atención de plazos, además
de selección de autores, reasignación y restauración de capturas y artefactos.
Las pruebas nuevas contrastaron los extremos inclusivo y exclusivo del periodo,
varios acuerdos dentro de una sesión y varias medidas dentro de una decisión.
Correcciones, retiros administrativos y reintentos no añadieron actuaciones;
una contribución original permaneció atribuida a su cuenta aunque terminara su
asignación. Un autor de otro expediente no quedó disponible para el solicitante.

La comprobación de audiencias respetó el conflicto de operación ya existente:
la recuperación de la revisión original conservó el recibo exacto. La decisión
cautelar repetida conservó su resultado original, conforme a su propio contrato.
No se alteraron estos flujos para adaptar sus recibos al informe. La restauración
probó las capturas y archivos existentes; la incorporación de fuentes nuevas no
recalcula una captura previamente guardada. Esta campaña focal no sustituye la
regresión global, la aceptación HTTP completa ni la observación de Qadra.

El recorrido ampliado con servidor real obtuvo un documento, nueve actuaciones
y un plazo atendido, con identidad, periodo y digest coincidentes en PDF y CSV.
Las actuaciones fueron una resolución, una notificación, los cinco tipos de
actos de recursos, un resultado de audiencia y una decisión cautelar. La sesión
contenía dos acuerdos y la decisión, dos medidas; sus contenidos no alteraron
el conteo esperado. La prueba confirmó los rechazos o recibos de los reintentos
según cada contrato y recuperó las revisiones originales. El navegador informó
26.8 segundos; la campaña completa, incluida preparación y compilación, duró
279.147 segundos, con un único worker y temporales en disco.

La evidencia previa de reasignación conserva un recorrido diferente: un autor
cargó un documento, perdió su asignación y fue seleccionado por otro litigante
con acceso vigente. Su informe de un documento, cero actuaciones y cero plazos
conservó PDF, CSV y aviso tras nueva autenticación. Esa ejecución duró
14.3 segundos de navegador y 330.290 segundos de campaña. No se suman ambas
observaciones como si fueran una única ejecución ni una evaluación de usabilidad.

La campaña compuesta de \texttt{api-demo.sh} aprobó en 562.754 segundos con
PostgreSQL 16.15 y servicios privados. Conservó el flujo HTTP de informes de
estado, sus archivos y avisos, y las comprobaciones de acceso después de
reiniciar y restaurar. Esta aceptación verifica la compatibilidad del servidor
compuesto; las cantidades positivas del catálogo de actividad corresponden al
recorrido de Qadra descrito anteriormente.
